


\section{Conclusion}
Now that we have landed on a well-performing social choice function for our coalition $\C$, we can also understand the associated deterrence behavior for the signaling game in \cref{sec:the-model}. 

If we omit the unlikely scenario that the probability vectors are almost random, $p=[0.55,0.55,0.55,0.55]$ and $q=[0.45, 0.45, 0.45,0.45]$, we obtain an average retaliation probability 
\begin{align*}
    \overline{R(\tau^*)}
    &= \frac{1}{13}\sum_{k=1}^{13}\Pr\left(\sum_{i=1}^4 v_i \geq 2 \mid T=1, p = p_k, q=q_k\right) \\
    &\approx 0.92
\end{align*}
where $p$ and $q$ ranges over the remaining $13$ probability vectors in \cref{table:p-q-vectors}. Similarly, 
\begin{align*}
    \overline{F(\tau^*)}
    &= \frac{1}{13}\sum_{k=1}^{13}\Pr\left(\sum_{i=1}^4 v_i \geq 2 \mid T=0, p = p_k, q=q_k\right) \\
    &\approx 0.12.
\end{align*}
On average, across a plethora of scenarios, the adversary state $\A$ attacks when the perceived benefit and costs satisfy 
\[0.92 \leq \frac{B_1}{B_1+C_1},\]
which happens if the costs are small compared to the benefits. Hence, for four-member coalitions with sufficient deterrence technologies, like strategic nuclear warheads, deterrence will be quite credible in most situations when using the majority rule as their institutional design.

